\section{Conclusion}
To conclude, we proved completeness for all controlled $n$-partite states, which we showed form a commutative ring isomorphic to multilinear polynomials. Also, we showed that all controlled $n$-qubit square matrices form a non-commutative ring. Furthermore, we have completeness for plugging controlled states into the control wires of controlled diagrams, isomorphic to all multivariate polynomials over same-size square matrices, with application to factoring Hamiltonians. When the controls target mutually exclusive sectors, a rewrite rule can be applied to copy any controlled diagram, and thereby factor any Hamiltonian.

We have shown that every (controlled) state computes a polynomial; hence, we can interpret a universal fragment of the ZXW-calculus as corresponding to arithmetic circuits. This generalises Ref.~\cite{carette2023compositionality}, which found an algebraic interpretation of a certain fragment of ZW calculus.
This line of reinterpreting quantum circuits as computing polynomials rather than unitary matrices offers a new perspective on quantum computation.

In another direction, we can apply completeness for polynomials isomorphic to controlled states to study entanglement. It can be easily shown diagrammatically that the polynomials corresponding to entangled (non-separable) states are exactly those that cannot be factored into irreducibles containing only variables corresponding to Alice's subsystem or only corresponding to Bob's subsystem. Since there are efficient algorithms for polynomial factorisation ~\cite{forbes2015complexity}, this gives rise to a novel entanglement classification algorithm for pure states. Further developing this into a more refined algebraic theory of entanglement, building on the work in Ref.~\cite{Agnew2023Masters}, could offer further insights.

We expect these results admit a natural extension for any controlled square matrices and controlled states in finite-dimensional Hilbert space, due to the completeness proof techniques of~\cite{poor2023completeness,deVisme2025minzw, Poor2025ZXFD}.
In all three of these works, while the central completeness result was for linear maps of higher- or mixed-dimensional Hilbert spaces, the normal forms themselves all involved restricting a qudit input to the W node to a qubit subspace of the qudit. Given the similarity of the ZW rules, we expect the underlying ring structures to be isomorphic even though \cite{poor2023completeness,Poor2025ZXFD} and~\cite{deVisme2025minzw} differ in interpretation of the generators, where the W node of the latter is rescaled by multinomial coefficients relative to the former.
Note that the completeness results of this work for controlled states do not directly follow from these completeness results for states, whose normal forms fix the control wire of the topmost W node to $\ket{1}$ instead of leaving this control wire open-ended.

As for when the control wires can take values beyond this two-dimensional subspace, we expect there to likewise be interesting underlying algebraic structure.
A starting guess would be that if qudit controlled states also admitted an equivalence to multivariate polynomials, they would instead be modulo $d$th powers of the variables, of the form $\mathbb{C}^{d-1}[x_1,...,x_n]/({x_1}^d,...,{x_n}^d)$, due to the Hopf law between Z and W.
Qudit multiple-control would likely have more complex structure than the qubit case here, considering the constructions for all prime-dimensional $d$-ary classical reversible gates and W states built in Refs.~\cite{Yeh2022QutritCtrl, Roy2023quditzh, Yeh2023quditW}.

We are broadly interested in exploring the relations between controlled diagrams, controlled quantum circuits, and Hamiltonian operators.
The completeness result for controllable quantum circuits established in~\cite{delorme2026ctrl} circumvents the unboundedness issues inherent to quantum circuits defined as a prop~\cite{clement2024minqc}, indicating that control as a higher-order operation on quantum circuits is a powerful construction for equational reasoning. In this work, we rewrite controlled diagrams to their fully expanded polynormal form, therefore establishing that equational reasoning on the quantum circuits as diagrammatical objects~\cite{shaikh2022sum} suffices to achieve any possible factorisations of Hamiltonian operators.
While the completeness of the rules means these rules suffice to derive any equalities of quantum circuits, more work can be done on how to effectively apply these rules towards the objective of synthesising more compact circuit decompositions, particularly for multiple-controlled gates and for block encodings.

Moreover, we are curious about how to embed these new semantics for quantum controlled states and matrices into a functional programming language like in Ref.~\cite{rennela2020clctrllinlogic}, and to interpret our diagrams in mixed-state quantum mechanics and relate them to the operational and denotational semantics for coherent control of Ref.~\cite{barsse2025qplcoherent} and to the channel construction of Ref.~\cite{deFelice2026dataflow}.
One approach is to consider using sector-preserving channels~\cite{Vanrietvelde2021ctrlsector} and scoped effects~\cite{lindley2024scoped} to more generally formulate semantics of multiple-control for heterogenous Hilbert spaces.

Yet another direction is to reconcile the interpretation of diagrammatic differentiation of our arithmetic polynomial circuits by the approach in Ref.~\cite{wilson2023diffpolycirc}, with that of quantum circuits and ZX diagrams in Refs.~\cite{toumi2021diagdiff, Wang2024diffintzx, jeandel2024adddiffzx}.